\section{Reusable proofs and source boundaries}

The mathematical chapters follow finite counts and the estimates they need.
The source tree also factors recurring technical arguments into reusable
interfaces. The following examples are useful when reading an individual proof:
\begin{itemize}
\item A finite prime-pair transport theorem reindexes a supported integer
      sum by its unique admissible ordered prime pair. Concrete cutoffs,
      filters and kernels remain inputs to that theorem. It lives with
      the Liu weights in \texttt{MathlibNt.SieveTheory.Liu.Weights.LiuWeight}.
\item Uniform logarithmic error absorption separates the elementary
      eventual estimate from the singular-series specialization. The
      generic argument lives in \texttt{MathlibNt.Analysis.LogScaleAbsorption}.
\item The integral excess-cover theorem combines a local approximation on
      a common region with a finite cover of its boundary excess. It lives
      in \texttt{MathlibNt.Analysis.IntegralExcessCover}; each application
      supplies its actual measure, strips and logarithmic kernel.
\item Filtered output sums retain every label in each fibre, with arbitrary
      additive weights. \texttt{MathlibNt.SieveTheory.FiniteFibres} shares this
      transport across the B8, B10, S2, G11 and G12 source models. Unlike the
      injective prime-pair transport, it permits repeated outputs. The
      integer and real product-grouping identities use Mathlib directly.
\item Labelled incidence bounds in
      \texttt{MathlibNt.SieveTheory.FiniteLabelCounting} count reverse
      relation fibres, retaining repeated products as separate labels.
      Its fixed-output sigma-fibre lemma keeps both positivity and the
      product upper bound needed to cancel natural-number subtraction.
\item \texttt{MathlibNt.SieveTheory.FiniteRealWindows} writes an arithmetic
      window $(L,U]$ as a difference of closed prefixes. The predicate
      retains primality and the product congruence in the G11/G12 uses;
      $0\le L\le U\le N$ is supplied by the original geometric bounds.
      These exact AP-window counts are separate library interfaces, not
      dependencies of the current headline counting proof.
\item \texttt{MathlibNt.Analysis.MovingIntervalIntegral} shares indicator-section
      and Fubini arguments for B8/B9/B10, the Liu source triangle, the
      weighted Fouvry G9 region and B9-high. For a measurable outer set
      $s$ and measurable region
      $R=\{(u,v):u\in s,\ l(u)<v\le r(u)\}$, an integrable real kernel on $R$
      satisfies $\int_R f=\int_{u\in s}\int_{v\in(l(u),r(u)]}f(u,v)$.
      No endpoint ordering or continuity is added; empty sections remain
      empty. Each application still proves its own kernel integrability
      and its conversion from oriented interval integrals. This continuous
      identity makes no assertion about discrete prime boundary atoms.
\item \texttt{MathlibNt.Analysis.LogGridEstimates} shares reciprocal-kernel
      variation, exact half-open cell disjointness, weighted cell evaluation
      and strip integration across Liu and B8/B9/B10. The true logarithmic
      density and the application-specific excess cover are retained. Its
      three-strip bound is one-sided; Liu and B10 still supply separate
      lower bounds before concluding convergence.
\item The boundary regularity development first proves joint continuity of
      the complete moving integral under a fixed local integrable majorant.
      Compactness supplies the original uniform moduli without repeatedly
      splitting common intervals and short boundary strips. Positive lower
      screens, the original depth induction and the zero-depth jump sets
      remain; no explicit rate follows from this compactness argument.
\item The Jurkat--Richert auxiliary estimates bound a smooth reciprocal tail
      by a finite weighted moment and its Euler-product bound, taking the
      finite upper endpoint to infinity with all other parameters fixed.
      A shared normalization of the existing Selberg estimate at the (4.1)
      cutoff keeps its complete squared-count error in both ratio branches.
      The original cutoff definitions, constants and threshold order remain.
      The auxiliary bounds now share these finite-moment and normalization steps.
\item The damped Perron estimates share a two-majorant integral bound and a
      single-character sharp-to-smoothed comparison. With natural $M\ge3$, rectangular
      coordinates in $[1,M]$, $Q>0$ and $S\subseteq[1,Q]$, smoothing at
      $\epsilon=M^{-2}$ gives the factors $1/2+(7\log M+2)/\pi$ for a common
      real smoothing parameter in $[1/2,M+1/2]$ and $1/2+(14\log M+4)/\pi$ for
      natural-valued selectors $Y(q,\chi)\le M$. These multiply the same
      aggregate rank-one large-sieve bound; the selector keeps four rank-one
      contributions. Applying the sharp comparison at each cutoff and summing
      with the nonnegative weight $q/\varphi(q)$ adds $8/(\pi M)$ times the
      product of the rectangular coefficient $L^1$ mass and weighted family mass. The
      single-character comparison uses natural $M\ge1$, natural $Y\le M$ and positive
      products; the aggregate sharp interfaces retain $mn\le M$.
\item \texttt{MathlibNt.Analysis.SieveNormalization} isolates elementary
      lower-bound multiplication and upper-bound exponential cancellation.
      The four Li--Liu applications retain their actual counts, main masses,
      singular series, remainder budgets and parameter ranges; the alpha
      case reuses the existing inverse-log absorption estimate.
\item \texttt{MathlibNt.Analysis.RealLogPowerThreshold} packages eventual
      domination of a fixed multiple of a real log power by a positive power.
      The Suzuki applications preserve thresholds before the varying sieve,
      the same-constant contract and the moving-depth domain. The remaining
      local algebraic compression is not a new analytic estimate.
\item \texttt{MathlibNt.Tactic.PolynomialDeriv} and
      \texttt{MathlibNt.Tactic.ElementaryDeriv} replace repeated derivative-rule
      assembly by kernel-checked proof generation. Polynomial formulas reuse
      Mathlib's polynomial derivative theorem; logarithmic and rational formulas
      retain their actual nonzero conditions and final algebraic identities.
      These tools change proof maintenance, not the sieve constants or ranges.
\end{itemize}
These interfaces preserve the counting objects, factor-size constraints,
signed sums and endpoint conventions. Both current quantitative entry points
consume the log-grid layer. The rewritten boundary-regularity and Rankin
proof bodies are used through their library interfaces, not by either
current headline quantitative proof.
The shared Perron and Suzuki arguments occur in both quantitative chains;
these four count normalizations occur in the Li--Liu chain. The formal proofs
share these analytic and algebraic steps.

The prime-number-theorem implementation shares the three-factor norm estimate
and the right-contour power identity. It recombines the nine contour pieces
with their original signs before taking norms, and separately records the
smoothing, outer, shifted and central error bounds. In particular, the two
infinite right-contour tails have scale $X\log X/(\epsilon T)$.

For a real interval $[z_1,z_2)$ with $2\le z_1\le z_2$, the inverse-product
argument takes $k=\max(2,\lceil z_1\rceil-1)$ and
$n=\max(k,\lfloor z_2\rfloor)$. The same integer comparison applies to empty
prime sets and small real endpoints. The resulting constant is chosen before
the finite prime set and the two real endpoints.

The finite-product bridge writes the actual sieve product as $2S_mP_m$,
where $m=Z-1$, $S_m$ is Liu's odd-prime truncation and $P_m$ is the ordinary
prime product up to $m$. Multiplication by $\log m$ gives both signs of the
Mertens estimate in one place. The lower and upper consumers retain their
own finite-to-infinite truncation comparisons and original parameter domains.

The Siegel--Walfisz estimates share a horizontal Mellin--Bochner bound,
the signed three-edge contour identity with two tails, and scalar tail and
smoothing estimates. Each application supplies its zero-free region and
logarithmic-derivative bound. The quadratic argument also retains the base
term used when transporting from a negative left exponent.

For small cutoffs, the Bombieri--Vinogradov argument bounds the actual
weighted coefficient energy, then the complete squared-error expression,
and finally absorbs logarithmic powers. The balanced and adaptive cases use the same
estimate with their actual cutoff. The adaptive cutoff is eventually bounded
by the balanced cutoff; monotonicity is not required. The modulus weight,
prefix maxima and harmonic factors stay in the expression being estimated.

The \sourcefile{MathlibNt/AnalyticNumberTheory/LargeSieve/StandardBVHighChosenUnconditional.lean}{high-conductor estimates}
reuse the quadratic Abel/conductor envelope from the
\sourcefile{MathlibNt/AnalyticNumberTheory/LargeSieve/StandardBVBlockL1WeightedPrimitive.lean}{block first-moment module}.
Writing $A_N$ for the Abel amplifier prefix maximum and $H_Q$ for the
conductor harmonic factor, it gives $4A_NH_Q^2\le48(\log N)^2$ for natural
$B$, $N\ge3$ and the actual Pan cutoff $Q$ at exponent $B$,
including $Q=0$. The chosen Type-I and Type-II estimates use this bound
directly; bare block sources keep the five-log reserve through its existing
corollary. The block assembler and the chosen small-term estimate also share
the square-to-first-moment step: at conductor threshold at least one,
$P^2(3HL)\le Z^2$ with $Z\ge0$ gives $P\,\mathrm{mean}\le Z$, where $H$ is the
high-conductor harmonic factor and $L$ is the actual primitive prefix-square
ledger on the same high-conductor set. The factor three stays in the payment.
This square comparison permits signed $P$; the first-moment block estimates
use a nonnegative multiplier.

For each character, the library's
\sourcefile{MathlibNt/AnalyticNumberTheory/LargeSieve/ConductorChangeLevelLedger.lean}{conductor-change ledger}
bounds its prefix maximum square by twice its primitive-conductor prefix
maximum square plus $8N$ times the coefficient energy on integers in $(M,M+N]$ not
coprime to the original modulus. It then sums with $q/\varphi(q)$ and regroups
the primitive term by conductor, avoiding a repeated error estimate under
nested sums. The principal term stays in its separate exact splitting identity.

The library's \sourcefile{MathlibNt/AnalyticNumberTheory/LargeSieve/ImprimitiveConductorWeightLinear.lean}{linear conductor transport}
applies the existing theorem for an arbitrary nonnegative family $F(d,\psi)$
directly to the primitive prefix maximum square. For natural $D>0$, the
imprimitive-weighted sum on $[D,2D]$ is bounded by the primitive-weighted sum
on $[1,2D]$ times $(Q/D)H_{Q/D}$, with natural-number division and the same
harmonic factor as above. This specialization supplies nonnegativity and
avoids repeating the weight comparison and conductor-window enlargement.

The Suzuki Case-I endpoint core uses the application's envelope-transport
inequality at a fixed sieve level to normalize the endpoint estimate. It then
combines three source-order remainders against the same nonnegative budget. Its strict and even-endpoint consumers retain separate
domains, with thresholds chosen before the varying sieve.

Write $\theta(x)=\sum_{p\le x}\log p$, summing over primes, and let $\pi(x)$
count primes at most $x$.
In the \sourcefile{PrimeNumberTheoremAnd/Consequences.lean}{prime-counting consequences}, Abel summation and the
Chebyshev remainder estimate give $\pi(x)\sim x/\log x$ from
$\theta(x)\sim x$. The logarithmic integral has the same scale: on $t>1$,
\[
  \left(\frac{t}{\log t}\right)'
    =\frac1{\log t}-\frac1{(\log t)^2},
\]
so the fundamental theorem of calculus and the squared-log integral bound give
$\int_2^x dt/\log t\sim x/\log x$, and hence
$\pi(x)\sim\int_2^x dt/\log t$. The endpoint identity includes $2\le a\le b$
with $a=b$ allowed. The Blueprint's epsilon-and-cutoff presentation gives
another proof of this asymptotic comparison; the Lean implementation chains
these equivalences. The legacy matcher interface is maintained separately.

\section{Finding and checking a declaration}

Each named result in the Blueprint links to a concrete declaration in the
source revision recorded by the website. The Lean documentation provides
its full type, surrounding definitions and module imports. For a major
estimate, read both the declaration and the definitions in its input type:
a packaged input can specify a source, density, modulus weight, cutoff and
order of quantifiers that a short theorem title cannot display.

The public entries are \texttt{Goldbach} for Chen,
\texttt{Goldbach.OnePlusOneNine} for Li--Liu, and \texttt{Goldbach.All} for
both. These entries expose independent theorem contracts over shared
foundations. Reusing a library does not identify different sieve densities
or different representation predicates.

Lean checks the formal proofs using propositional extensionality, classical
choice and quotient soundness as its standard logical foundations. The
project's acceptance process checks statements, compiles with warnings as
errors, inspects theorem axiom dependencies and independently replays
compiled declarations. The Blueprint has additional editorial obligations:
its mathematical statements and explanations must match the implementation,
and its reading route must expose the substantive estimates needed by each
conclusion. Graph extraction and rendering tests address a separate,
structural part of that verification.

\section{Mathematical sources and attribution}

This development builds on UyNewNas/chen-theorem-lean, the attributed
analytic-number-theory library, Mathlib and adapted
PrimeNumberTheoremAnd sources. Upstream copyright notices are retained
under Apache-2.0. The repository's \sourcefile{docs/PROVENANCE.md}{mathematical provenance guide} identifies the mathematical sources and their roles. The two proof chapters describe the
implemented arguments: classical results and modern replacements are
identified at the places where they enter the proof.

The website build record distinguishes its own editorial source revision,
the Blueprint revision and the Lean API revision. A local documentation
preview can reuse a previously generated API artifact while preserving
that artifact's original source links. Online deployment is a separate
step from generation and local verification.

\url{https://github.com/subfish-zhou/goldbach-lean}


\paragraph{Character estimates and shared payments.}
The reduced-residue character estimate specializes the existing Parseval identity on units. The primitive-character estimate retains its separate direct orthogonal-family Bessel argument, using the existing inner-product and energy identifications. These are distinct proof routes. In the pointwise Siegel--Walfisz estimates, the horizontal contour edges are controlled by the tail term. The shared masked dyadic payment for the Fouvry estimates belongs to the existing MaskedW module; the callers retain their geometric hypotheses and threshold choices.


\paragraph{S2 limit budget and Vaughan cell estimate.}
The S2 main-mass estimate combines the weighted logarithmic-kernel limit with its vanishing remainder contribution before choosing the final threshold. The comparison still uses the actual prime carrier, a larger prime window with nonnegative weights, and the original logarithmic coefficient. For the Vaughan all-aspect cell estimate, one local symmetric square-root estimate controls both cross terms. The argument retains the actual conductor cell, active rectangles, normalization, and the original four-term bound; the symmetric estimate is local to the proof rather than a new public interface.


\paragraph{Weighted Perron norms and the S1 error budget.}
The selector estimate first bounds the nonnegative weighted sum of norms of the damped Perron integrals. The four selector-independent phase twists are controlled by the rank-one large-sieve energy bridge, and the two integrable majorants preserve the exact coefficient 14 log M + 4 before the signed Perron identity is averaged. The fixed-s S1 bound keeps the true carrier, sieve product and singular-series normalization: density and main-mass losses share the scalar delta budget, while half of that budget absorbs the paid Bombieri--Vinogradov remainder. A common threshold combines these estimates; S1 depends transitively on the selector rather than forming an independent branch.
